% Transcription — fonds Alexandre Grothendieck, shelfmark 122, batch 1 (pages 1-20).
% Established from the facsimile archives/batches/122/batch-01.pdf (prepared
% from a copy of 122.pdf supplied by the user; PDF page k = archivists' page
% k, no cover sheet), per the skill transcribe-grothendieck. The folder is
% thirty-four pages; this is the first of two batches. Batch 2 (pages 21-34)
% was transcribed first, from its own pages alone; its header and notation
% were read before this pass and are followed here.
%
% Pass: Opus 5.5 (claude-opus-5-5), 2026-09-24 — first pass, unchecked against the pages by a human.
%
% Pages skipped: 1 and 18, blank cream sheets; 3, 5, 7, 9, 11, 13, 15, typed
% leaves of a Bourbaki draft (« n° 387 », pp. -42- to -48-, exercises on
% valuations, linearly compact modules and the « dual torique ») with nothing
% of his; 17, a typed announcement from the Institut Henri Poincaré (lectures
% by N. Efimov, first lecture 25 April 1968), nothing of his; 20, a leaf
% (-4.15-) of the English typescript on geometric quotients whose other
% leaves are the versos of batch 2. All skipped by the settled rule, without
% description in the body. No page is summarised.
%
% Pages transcribed: 2, 4, 6, 8, 10, 12, 14 (his ink on tan sheets, the backs
% of the Bourbaki leaves), 16 (a sheet of sketches), 19 (white paper, blue-black ink, pencil circles —
% the same paper and ink as pp. 21-31 of batch 2). No pagination of his is
% visible.
%
% What the batch is about. Topological preliminaries to batch 2. Page 2: P a
% (para)compact manifold with boundary of pure dimension n, P̊ = P − ∂P → P a
% homotopy equivalence, H_∂P(P,F) = 0 for F locally constant; the exact
% sequence H_!(P̊) → H(P) → H(∂P) with H_!^i(P̊,F) ≃ H^{n-i}(P,F̌⊗T)^∨, hence
% φ ∈ H^i(P,F) ⊗ H^{n-i}(P,F̌⊗T), Poincaré duality when ∂P = ∅; a collar
% ∂̃P ≃ ∂P × D^1 and P' = closure of P − ∂̃P. Page 4 (a block cancelled by
% long diagonals, then a clean continuation): H^i_{P'}(P,F) ≃ H^i_!(P̊,F) by
% excision (P',∂P') → (P,P−P'). Page 6: W = P ∪ Q with P ∩ Q = ∂P = ∂Q = R,
% W − R = P̊ ⊔ Q̊, the relative sequence for (W,R,W−R) read as Mayer-Vietoris,
% and its compactly supported « transposée » through Poincaré duality. Page 8:
% the relative sequence of P ⊂ W, H^i_P(W,F) ≃ H^i_!(P̊,F) by excision on a
% small thickening P̃, hence H_!(P̊) → H(W) → H(Q). Page 10: a hexagon of the
% four sequences around H^i(W,F), a boxed table of the sequences S(W,P),
% S_!(W,Q), S(W,R) = MV(W;P,Q) and their duals, a cancelled block, and a boxed
% « Récapitulation » Ψ, Φ, Φ', Ψ' with (Φ,Φ') and (Ψ,Ψ') dual. Page 12: W
% glued from two fibre bundles P → S (fibre M) and Q → T (fibre N) with
% α : S ≃ ∂N, β : T ≃ ∂M and trivialisations σ, τ of the boundaries, giving
% φ : ∂P ≃ ∂Q. Page 14: W obtained by gluing along φ; W − Q = P − ∂P ↪ P and
% W − P = Q − ∂Q ↪ Q homotopy equivalences; a diagram of the pieces (N a disc,
% ∂N = S^1). Page 16: sketches — the dictionary W, P, Q, R ↔ V', X'_η = V_η,
% V'_s, V'_s × η and M, N, T, S ↔ X'_η, e, V'_s, η ~ S^1 (so the manifolds
% above model a Milnor-type fibration); two cube diagrams; a count of 40
% (« en fait 80 ») exact sequences. Page 19: the four triangles (a), (b), (c),
% (d) in the derived category of Λ[π] that batch 2 invokes: (a) K^• =
% R(Γ,π)(M) → L^{•triv} = RΓ(W_0)^{triv} (ρ, compatible with cup) with third
% vertex R(Γ,π)_!(M̊); (b) its RΓ^π, with (α,β) : RΓ(P) = RΓ(W−W_0) → RΓ(W_0) ×
% RΓ(W_0)(−1); (c) its forgetful image; (d) the triangle RΓ(W_0)[−2] →
% RΓ_!(W) → RΓ(W−W_0) with « le β de (b) ».
%
% Notation fixed once, following batch 2: \mathbb{R}\Gamma, \mathbb{R}(\Gamma,
% \pi), \mathbb{R}\Gamma_{!}; K^\bullet, L^\bullet; \mathring{P}, \mathring{M}
% for ring-topped letters; (-)^\vee for the dual; the circled letters naming
% triangles as (a), (b), (c), (d). His orientation sheaf, a script letter
% like a C with a tail, is given as \mathcal{T}_W (plain T on p. 2). His ∂ and
% δ for the boundary are both given as \partial; the thickened boundary as
% \widetilde{\partial P}.
%
% Check against batch 2: the exponent of K^• in (c) is, on p. 19, a small
% « ω » (or « w »), K^{•ω}, R^ω; batch 2 read the same exponent in (c') as
% « ens » (doubtful). Not changed there; flagged in the report.
%
% Hardest pages and readings to check first: p. 4 (cancelled block, the
% « mod » readings, the subscripts of H^i_{P'}); p. 10 (orientation, the
% cancelled Mayer-Vietoris block, the arrowheads of the hexagon); p. 14 (the
% arrowheads of the diagram); p. 16 (cube diagrams only described); the
% connective prose throughout.

\documentclass[11pt,a4paper]{article}
% Le préambule commun à toutes les transcriptions du fonds Grothendieck.
%
% Il tient en une idée : **l'incertitude se compose, elle ne se résout pas.**
% Une transcription qui lisse un mot douteux a détruit la seule chose pour
% laquelle on la faisait — un lecteur doit pouvoir savoir, à chaque endroit,
% ce qui était sur le feuillet et ce qui a été ajouté. Les six macros qui
% suivent sont donc l'appareil critique, et non de la décoration.
%
% Les mêmes commandes sont reconnues par `scripts/render.mjs`, qui produit la
% vue de lecture du site. Une macro ajoutée ici doit l'être aussi là-bas,
% faute de quoi le rendu échouera bruyamment — ce qui est voulu : un rendu
% silencieusement faux est pire qu'un rendu absent.

\ProvidesPackage{grothendieck}[2026/08/08 Transcriptions du fonds Grothendieck]

\usepackage{amsmath,amssymb,amsthm}
\usepackage{mathrsfs}
\usepackage{stmaryrd}
% Les diagrammes commutatifs sont partout dans ce fonds : c'est la langue
% même de la géométrie algébrique de Grothendieck, et une transcription qui
% les rendrait en prose ne transcrirait rien.
\usepackage{tikz-cd}
% Français par défaut — c'est la langue des feuillets — mais l'anglais est
% chargé aussi : la traduction et le résumé sont des documents anglais, et la
% typographie française (espaces avant les deux-points) y serait une faute.
% Ces documents basculent par \selectlanguage{english} en tête de corps.
\usepackage[english,french]{babel}
\usepackage{fontspec}
\usepackage{geometry}
\geometry{a4paper,margin=25mm,marginparwidth=28mm}
\usepackage{marginnote}
\usepackage{xcolor}
% ulem, not soul. `soul`'s \st and \ul are built on a character-by-character
% scan that XeLaTeX's font machinery breaks: the document fails with an
% undefined control sequence rather than degrading. ulem does the same two jobs
% and is Unicode-engine safe. `normalem` stops it hijacking \emph, which the
% transcriptions use for Grothendieck's own emphasis.
\usepackage[normalem]{ulem}
\usepackage{hyperref}
\hypersetup{colorlinks=true,linkcolor=black,urlcolor=black,pdfborder={0 0 0}}

% --- Métadonnées du lot -----------------------------------------------------
% Reprises telles quelles par le script de rendu, qui les affiche en tête de
% la vue de lecture. Elles fixent ce que le fichier prétend couvrir : sans
% elles, une transcription flotte sans point d'attache dans un fonds de seize
% mille feuillets.

\newcommand{\folder}[1]{\gdef\grothFolder{#1}}
\newcommand{\batch}[1]{\gdef\grothBatch{#1}}
\newcommand{\leaves}[2]{\gdef\grothFirst{#1}\gdef\grothLast{#2}}
\newcommand{\foldertitle}[1]{\gdef\grothFoldertitle{#1}}
\gdef\grothFolder{?}\gdef\grothBatch{?}\gdef\grothFirst{?}\gdef\grothLast{?}\gdef\grothFoldertitle{}

% --- Le résumé --------------------------------------------------------------
%
% Il ouvre la lecture modernisée, sous le titre « Résumé », et il est composé
% en petit. Ce n'est pas un ornement : c'est le seul passage du document qui
% ne soit pas des mathématiques, et un lecteur doit pouvoir le reconnaître
% avant de l'avoir lu — puis le sauter s'il connaît déjà le sujet, ou s'y
% arrêter s'il ne le connaît pas. Le filet qui le suit marque la reprise du
% corps.

\newenvironment{resume}
  {\small\setlength{\parskip}{0.45em}}
  {\par\medskip\hrule\medskip}

% --- Mots-clés ---------------------------------------------------------------
%
% En anglais, en clôture du résumé de la lecture modernisée : le vocabulaire
% moderne sous lequel chercher ce que les feuillets construisent. Le manifeste
% du site (`npm run manifest`) les extrait pour étiqueter la cote entière —
% cette ligne est la source unique des tags, il n'y a pas de fichier à côté.
\newcommand{\keywords}[1]{\par\smallskip\noindent
  {\footnotesize\textsc{Keywords} --- \textit{#1}}\par}

% --- L'appareil critique ----------------------------------------------------

% Le numéro de feuillet, en marge. C'est l'ancre qui permet de régler un
% désaccord sur une formule en regardant *un* feuillet plutôt qu'une chemise
% de six cents. Le site s'en sert aussi pour tourner les pages du fac-similé
% quand on fait défiler la transcription.
\newcommand{\leaf}[1]{%
  \marginnote{\footnotesize\color{gray!70}\textsf{#1}}%
  \label{leaf:#1}\ignorespaces}

% Illisible : on ne devine pas. Un mot inventé qui se lit comme les autres est
% le pire résultat possible.
\newcommand{\ill}{\textcolor{red!60!black}{[\,\ldots\,]}}

% Lecture douteuse : le mot est donné, et signalé comme incertain.
\newcommand{\uncertain}[1]{\uline{#1}}

% Ajout de l'éditeur — une lettre manquante, un mot sous-entendu.
\newcommand{\add}[1]{\textcolor{blue!50!black}{[#1]}}

% Biffé par Grothendieck lui-même. Ce qu'il a rayé fait partie du document :
% une rature dit souvent où la pensée a changé de direction.
\newcommand{\struck}[1]{\sout{#1}}

% Note du transcripteur, sur l'état matériel du feuillet.
\newcommand{\note}[1]{\par\smallskip{\small\color{gray!60!black}\textit{[#1]}}\par\smallskip}

% Note marginale de Grothendieck, distincte du corps du texte.
\newcommand{\marginal}[1]{\par\smallskip\noindent\hspace{1em}%
  {\small\color{gray!60!black}#1}\par\smallskip}

% --- Titre ------------------------------------------------------------------

\AtBeginDocument{%
  \begin{center}
    {\large\bfseries Fonds Alexandre Grothendieck --- cote n\textsuperscript{o}~\grothFolder}\\[2pt]
    {\itshape\grothFoldertitle}\\[4pt]
    {\small feuillets \grothFirst--\grothLast\ (lot \grothBatch)}\\[2pt]
    {\footnotesize\color{gray!60!black}Transcription établie d'après le fac-similé numérisé
     par l'Université de Montpellier.\\
     Les lectures incertaines sont soulignées, les illisibles notées [\,\ldots\,],
     les ajouts entre crochets.}
  \end{center}
  \vspace{1em}}

\endinput


\folder{122}
\batch{1}
\pages{1}{20}
\dating{[vers 1968-1971]}
\watermark{Édition de démonstration}
\foldertitle{Point singulier isolé de la fibre : notes manuscrites (s.d.)}

\begin{document}

\page{2}

$P$ variété à bord \add{para}compacte, purement de dimension $n$
\note{« para » est ajouté au-dessus de la ligne et entouré}

$\partial P$ le bord

$\mathring{P} = P - \partial P$

$\mathring{P} \longrightarrow P$ est une équivalence d'homotopie, donc
\[
  H^{*}(P, F) \longrightarrow H^{*}(\mathring{P}, F) \quad \text{un isom}
  \qquad (\text{si } F \text{ loc. constant})
\]
\struck{donc} i.e.
\[
  H^{*}_{\partial P}(P, F) = 0
\]
et en localisant
\[
  \underline{H}^{*}_{\partial P}(F) = 0 .
\]
Considérons d'autre part la suite exacte
\[
  \cdots \longrightarrow H^{i}_{!}(\mathring{P}, F) \longrightarrow H^{i}(P, F)
  \longrightarrow H^{i}(\partial P, F) \longrightarrow \cdots
\]
\marginal{si on travaille avec faisceaux loc. constants \ill{}, i.e. \uncertain{fibres}
de dim finie}
\[
  H^{i}_{!}(\mathring{P}, F) \simeq H^{n-i}(\mathring{P}, \check{F} \otimes T)^\vee
  \simeq H^{n-i}(P, \check{F} \otimes T)^\vee
\]
\note{les deux isomorphismes sont écrits verticalement sous
$H^{i}_{!}(\mathring{P}, F)$ ; $T$ désigne ici le faisceau d'orientation}

Ainsi, on a des hom\add{omorphismes}
\[
  \varphi : H^{n-i}(P, \check{F} \otimes T)^\vee \longrightarrow H^{i}(P, F)
\]
i.e.
\[
  \varphi \in H^{i}(P, F) \otimes H^{n-i}(P, \check{F} \otimes T)
\]
[si $\partial P = \emptyset$, les $\varphi$ sont des isom, ce qui est la
dualité de Poincaré pour la variété \uncertain{topologique} $P$].

Soit d'autre part $P' \subset P$, $P' = \overline{P - \widetilde{\partial P}}$,
$\widetilde{\partial P}$ un « épaississement » de $\partial P$,
\[
  \widetilde{\partial P} \simeq \partial P \times D^{1} .
\]
\note{en bas à gauche, un croquis : un anneau hachuré $\widetilde{\partial P}$
bordé par le cercle $\partial P$}

On a donc \struck{\ill{} de toutes \uncertain{sortes}}

\struck{$\partial P \to \widetilde{\partial P} \leftarrow \partial P'$, $\partial P' \to P$
équivalences d'homotopie}\note{deux lignes biffées ; la seconde est surchargée}

\page{4}

\note{tout le haut de la page, jusqu'au trait horizontal, est barré de longs
traits obliques et entouré d'un grand trait courbe}

Dans \ill{}, l'inclusion $(P, \widetilde{\partial P}) \subset (P, \partial P)$
\ill{} un hom\add{omorphisme} de suites exactes
\[
  \begin{array}{ccccccc}
    \to & H^{i}_{!}(P - \widetilde{\partial P}) & \to & H^{i}(P) & \to & H^{i}(\widetilde{\partial P}) & \to \cdots \\
        & \downarrow & & \| & & \downarrow{\wr} & \\
    \to & H^{i}_{!}(P - \partial P) & \to & H^{i}(P) & \to & H^{i}(\partial P) & \to \cdots
  \end{array}
\]
d'où \ill{} \ill{} \ill{} \ill{}
\[
  H^{i}_{!}(P - \widetilde{\partial P}) \simeq H^{i}(P - \partial P)
\]
\note{à gauche et à droite, deux croquis de flèches entre
$\partial P$, $\widetilde{\partial P}$, $\partial P'$, $\widetilde{\partial P} - \partial P$,
$\widetilde{\partial P} - \partial P'$ et $P - P' = \widetilde{\partial P} - \partial P'$,
avec la mention « équivalences d'\ill{} »}

\[
  \begin{array}{ccccccc}
    \to & H^{i}_{P'}(P, F) & \Longrightarrow & H^{i}(P, F) & \to & H^{i}(P - P', F) & \to \cdots \\
        & \downarrow & & \| & & \updownarrow{\scriptstyle 2} & \\
    \to & H^{i}(P, \partial P; F) & \longrightarrow & H^{i}(P, F) & \to & H^{i}(\partial P, F) & \to
  \end{array}
  \qquad
  \begin{array}{c} (P, P - P') \\ \uparrow \\ (P, \partial P) \end{array}
\]
\note{le premier terme de la seconde ligne peut se lire
$H^{i}(P/\partial P; F)$ ; à côté de la flèche verticale de droite, un
$H^{i}(\partial P)$ biffé}

On trouve
\[
  H^{i}_{P'}(P, F) \simeq H^{i}(P \bmod \partial P, F) \simeq H^{i}_{!}(\mathring{P}, F)
\]
\note{« mod » est une lecture douteuse}
\ill{} \ill{} \ill{}
\[
  H^{i}_{P'}(P, F) \xrightarrow{\ \sim\ } H^{i}_{!}(\mathring{P}' \bmod \partial P', F)
  \simeq H^{i}_{!}(\mathring{P}', F)
\]
\note{« mod » est ici biffé et surchargé ; l'indice du deuxième $H$ est douteux}
grâce à l'excision
\[
  (P', \partial P') \to (P, P - P') .
\]

\page{6}

$W$ variété \struck{\ill{}}

$P \subset W$ sous-variété à bord telle que $P = \overline{\mathring{P}}$.

Soit $Q = P - \mathring{P}$.\note{sic ; on attendrait $\overline{W - P}$}

Alors $Q$ est une sous-variété à bord satisfaisant $Q = \overline{\mathring{Q}}$, et
\struck{si}
\[
  W = P \cup Q, \qquad P \cap Q = \partial P = \partial Q .
\]
Si $W$ est purement de \uncertain{dimension} $n$, alors $P$ et $Q$ aussi, et
$P \cap Q = \partial P = \partial Q$ est purement de dim\add{ension} $n - 1$.

Soit $R = \partial P = \partial Q$. \ill{} \struck{\ill{}} \ill{} donc
\[
  W - R = \mathring{P} \amalg \mathring{Q} .
\]
\struck{$H^{i}(W) \to H^{i}(W)$}\note{ligne biffée, lecture douteuse}

\begin{tikzcd}[column sep=small, row sep=small]
 & & \mathring{Q} \arrow[d] \\
 & W & Q \arrow[l, hook'] \\
\mathring{P} \arrow[r, hook] & P \arrow[u, hook] & P \cap Q = R = \partial P \arrow[l, hook'] \arrow[u, hook]
\end{tikzcd}

\note{croquis dans la marge gauche ; sous $P \cap Q = R$, en colonne :
$\partial P = T$, le $R$ étant surchargé}

La suite exacte de coh\add{omologie} relative pour $W$, $R$, $W - R$ donne
\[
  \cdots \to H^{i}(W, F) \to H^{i}(W - R, F) \to H^{i+1}(R, F \otimes \mathcal{T}_{R})
  \to H^{i+1}(W, F)
\]
\[
  H^{i}(W - R) = H^{i}(\mathring{P}) \times H^{i}(\mathring{Q}) = H^{i}(P) \times H^{i}(Q)
\]
\marginal{faisceau d'orientation normal : $T$ dans $W$, ici $\simeq \mathbb{Z}_T$}\note{dans une
bulle reliée à $\mathcal{T}_R$ ; l'indice de $\mathcal{T}$ est surchargé, $R$ ou $T$}

Cette suite exacte relative n'est sans doute autre que la suite exacte de
Mayer-Vietoris associée à $(*)$, \ill{} aussi [\ill{} \uncertain{forme}
\uncertain{transposée}]
\[
  \cdots \to H^{i}_{!}(W - T, F) \to H^{i}_{!}(W, F) \to H^{i}_{!}(R, F) \to \cdots
\]
\[
  \begin{aligned}
  H^{i}_{!}(W - T, F) &\simeq H^{i}_{!}(\mathring{P}, F) \times H^{i}_{!}(\mathring{Q}, F) \\
  &\simeq H^{n-i}(\mathring{P}, \check{F} \otimes \mathcal{T}_W)^\vee \times
     H^{n-i}(\mathring{Q}, \check{F} \otimes \mathcal{T}_W)^\vee \\
  &\simeq H^{n-i}(P, \check{F} \otimes \mathcal{T}_W)^\vee \times
     H^{n-i}(Q, \check{F} \otimes \mathcal{T}_W)^\vee
  \end{aligned}
\]
\note{les isomorphismes sont écrits en colonne sous le premier terme. Dans la
dernière ligne, la lettre lue $Q$ ressemble au $R$ de la page ; le symbole lu
$\mathcal{T}_W$ (faisceau d'orientation) est une sorte de C à queue}

\page{8}

Écrivons aussi les 2 suites exactes relatives à l'inclusion $P \hookrightarrow W$. On a
\[
  \cdots \to H^{i}_{P}(W, F) \to H^{i}(W, F) \to H^{i}(W - P, F) \to \cdots
\]
\[
  H^{i}(W - P, F) \simeq H^{i}(\mathring{Q}, F) \xleftarrow{\ \sim\ } H^{i}(Q, F)
\]
\note{écrit en colonne sous $H^{i}(W - P, F)$. Sous $H^{i}_{P}(W, F)$, un
« $\simeq H^{i}_{!}(\mathring{P})$ » biffé, et dans la marge un mot biffé}

Or, si $\widetilde{P}$ est un petit épaississement de $P$ dans $W$, \ill{}
l'inclusion $P \to \widetilde{P}$ \ill{}, on \ill{} \ill{}
\[
  H^{i}_{P}(\widetilde{P}, F) \xrightarrow{\ \sim\ }
  H^{i}_{!}(\widetilde{P} \bmod \partial\widetilde{P}, F) \simeq H^{i}_{!}(\mathring{P}, F),
  \qquad
  H^{i}_{P}(\widetilde{P}, F) \simeq H^{i}_{P}(W, F) \ \ (\text{excision})
\]
\note{« mod » est une lecture douteuse ; les deux isomorphismes de droite
sont écrits verticalement. À gauche, un croquis : $P$ entouré d'un anneau
hachuré, l'épaississement $\widetilde{P}$}

Donc la suite exacte \uncertain{pour} l'inclusion $P \hookrightarrow W$ \uncertain{peut
s'écrire}
\[
  \cdots \to H^{i}_{!}(\mathring{P}, F) \to H^{i}(W, F) \to H^{i}(Q, F) \to \cdots
\]
\[
  H^{i}_{!}(\mathring{P}, F) \simeq H^{n-i}(P, \check{F} \otimes \mathcal{T}_W)^\vee
\]

\page{10}

\note{feuille tenue en largeur ; les parties sont écrites dans des sens
différents}

\begin{tikzcd}[column sep=tiny, row sep=small, nodes={font=\scriptsize}]
 & H^{i-1}(R, F) \arrow[dl, no head] \arrow[dr] \arrow[d] & \\
H^{n-i}(P, \check{F} \otimes \mathcal{T}_W)^\vee \simeq H^{i}_{!}(\mathring{P}, F) \arrow[r] & H^{i}(W, F) \arrow[dl] \arrow[dr] & H^{i}_{!}(\mathring{Q}, F) \simeq H^{n-i}(Q, \check{F} \otimes \mathcal{T}_W)^\vee \arrow[l] \\
H^{i}(P, F) \arrow[r] & H^{i}(R, F) & H^{i}(Q, F) \arrow[l]
\end{tikzcd}

\note{hexagone : les flèches de $H^{i}_{!}(\mathring{Q}, F)$ vers $H^{i}(W, F)$
et de $H^{i}(W, F)$ vers $H^{i}(P, F)$, $H^{i}(Q, F)$ sont repassées en gras ;
le trait de $H^{i-1}(R, F)$ vers le terme de gauche n'a pas de pointe
visible}

\marginal{suite exacte de cohomologie \ill{} pour $P \subset W$}
\marginal{suite exacte de cohomologie \ill{} $!$ pour $Q \subset W$}\note{dans
une bulle, de part et d'autre d'une flèche à deux pointes}

\[
  \boxed{\begin{array}{ccc}
    S_{!}(W, R) & \text{dual de} & S(W, R) \\
    \uparrow & & \downarrow \\
    S(W, P) \simeq S_{!}(W, Q) & \text{dual de} & S(W, Q) \simeq S_{!}(W, P) \\
    \uparrow & & \downarrow \\
    S(W, R) \simeq MV(W; P, Q) & \text{dual de} & S_{!}(W, R) \simeq MV_{!}(W, R)
  \end{array}}
\]
\note{le $P$ de $S(W, P)$ est surchargé ; les $R$ de la première ligne sont
douteux}

\struck{Suite exacte de Mayer-Vietoris}

\struck{Suite exacte de coh. relative \ill{}}
\[
  \begin{gathered}
  H^{i}_{Q}(W) \to H^{i}(W) \to H^{i}(W - Q) \simeq H^{i}(P) \\
  H^{i}_{Q}(W) \simeq H^{i}_{!}(W - P) = H^{i}_{!}(Q - \partial Q)
  \simeq \bigl(H^{n-i}(Q - \partial Q)\bigr)^\vee = H^{n-i}(Q)^\vee \\
  H^{i}_{Q}(W) \times H^{n-i}(Q)
  \end{gathered}
\]
\note{tout ce bloc, écrit dans l'autre sens, est barré de longs traits ; le
premier $\simeq$ sous $H^{i}_{Q}(W)$ est suivi d'un « ? ». En tête, à droite :
« $P$, $Q$, $P \cap Q = S \times T$, $W$ »}

\textbf{Récapitulation}

\[
  \boxed{\begin{array}{c|l}
    \begin{array}{ccc} & \Psi & \\ \swarrow & & \searrow \\ \Phi & & \Phi' \\ \searrow & & \swarrow \\ & \Psi' & \end{array}
    &
    \begin{array}{l}
      \Psi : H(P), H(Q), H(T), H(W) \\
      \Psi' : \ldots\ \text{dual} \\
      \Phi : H(P)^\vee, H(Q), H(W) \\
      \Phi' : H(P), H(Q)^\vee, H(W)
    \end{array}
  \end{array}}
\]
\note{dans la ligne $\Psi'$, un mot illisible devant « dual » ; après les lignes $\Phi$ et $\Phi'$, un mot illisible chacune, suivi
de « $H(T)$ »}
\[
  \left.\begin{aligned}
    \Phi &\simeq S(W, P) \simeq S_{!}(W, Q) \\
    \Phi' &\simeq S(W, Q) \simeq S_{!}(W, P)
  \end{aligned}\right\} (\Phi, \Phi' \text{ duals})
  \qquad
  \left.\begin{aligned}
    \Psi &\simeq S(W, T) \simeq MV(W; P, Q) \\
    \Psi' &\simeq S_{!}(W, T) \simeq MV_{!}(W; P, Q)
  \end{aligned}\right\} (\Psi, \Psi' \text{ duals})
\]

\page{12}

Soit $W$ une variété \struck{topologique}, diff\add{érentiable}.

$P$, $Q$ deux sous-variétés à bord

avec
\[
  P \cap Q = \partial P = \partial Q, \qquad P \cup Q = W
\]
et $P$ et $Q$ étant des fibrés diff.

\struck{Alors on}
\[
  \begin{array}{ccc}
    P \longrightarrow S & \qquad & Q \longrightarrow T \\
    \wr & & \uparrow \\
    M & & N
  \end{array}
\]
\note{dans la marge gauche, un $W$ biffé et « $\partial P$ »}

de bases $S$, $T$ des variétés diff\add{érentiables} compactes $S$, $T$, de
fibres des variétés à bords $M$, resp. $N$. Donc $\partial P$ est un fibré sur
$S$ de fibre $\partial M$, et $\partial Q$ un fibré sur $T$ de fibre
$\partial N$ :
\[
  \begin{array}{ccc}
    \partial P \longrightarrow S & \qquad & \partial Q \longrightarrow T \\
    \wr & & \uparrow \\
    \partial M & & \partial N
  \end{array}
\]
On suppose donnés des isomorphismes
\[
  \alpha : S \simeq \partial N \qquad \beta : T \simeq \partial M
\]
et \struck{que les fibrés} des trivialisations
\[
  \sigma : \partial P \simeq S \times \partial M \qquad
  \tau : \partial Q \simeq T \times \partial N
\]
d'où, par composition,
\[
  (\alpha, \sigma) : \partial P \simeq \partial N \times \partial M, \qquad
  (\beta, \tau) : \partial Q \simeq \partial M \times \partial N
\]
\note{le $N$ de $\partial N \times \partial M$ est surchargé}
et par suite un isom
\[
  \varphi : \partial P \simeq \partial Q .
\]

\page{14}

Ceci posé, on suppose $W$ construit en recollant $P$ et $Q$ suivant leurs bords
\struck{\ill{}} via $\varphi$.

\struck{On suppose} \struck{On aura}

L'hyp\add{othèse} \ill{} \uncertain{pour} $\partial P$ \uncertain{est} un fibré
trivial \uncertain{sur} \ill{} \uncertain{signifie} que l'inclusion \struck{\ill{}}
\[
  W - Q = P - \partial P \hookrightarrow P
\]
est une équivalence d'homotopie, \struck{\ill{}}

et de même pour
\[
  W - P = Q - \partial Q \hookrightarrow Q .
\]
\struck{\ill{} \ill{}}

Considérons le diagramme

\begin{tikzcd}[column sep=small, row sep=small, nodes={font=\scriptsize}]
 & & T = \partial M & & t \arrow[ll] \\
 & W & Q \arrow[l, hook'] \arrow[u, "2"'] & & N = D_{1} \sim e \arrow[ll, hook'] \arrow[u, no head] \\
\partial N = S^{1} & P \arrow[l] \arrow[u, leftrightarrow] & S^{1} \times T \arrow[l, hook'] \arrow[u, hook] & & \partial N = S^{1} \arrow[ll, leftrightarrow] \arrow[u, hook] \\
s \arrow[u] & M \arrow[l] \arrow[u, hook] & \partial M = T \arrow[l, hook'] \arrow[u, leftrightarrow] & & (s, t) \arrow[ll] \arrow[u, hook]
\end{tikzcd}

\note{le sommet central porte en entier
$S^{1} \times T = \partial N \times \partial M = \partial P = \partial Q = P \cap Q$ ;
à droite de ce sommet, un « $e$ » biffé. La flèche de $t$ vers $N$ est un
trait épais à deux extrémités marquées ; plusieurs têtes de flèche du
diagramme sont douteuses, reproduites telles que lues. « $D_{1} \sim e$ » est
une lecture douteuse}

\page{16}

\note{feuille de croquis, sans phrases suivies ; on en donne les inscriptions}

\[
  \begin{array}{cccc}
    W & P & Q & R \\
    V' & X'_{\eta} = V_{\eta} & V'_{s} & V'_{s} \times \eta
  \end{array}
  \qquad\qquad
  \begin{array}{ccc} & W & \\ P & & Q \\ & R & \end{array}
\]
\note{à droite, les quatre lettres disposées en losange et reliées par des
traits}

\[
  \begin{array}{cccc}
    M & N & T = \dot{M}, & S = \dot{N} \\
    X'_{\eta} & e & V'_{s} & \eta \sim S^{1}
  \end{array}
\]
\note{devant le $e$, un mot biffé ; le $N$ de $S = \dot{N}$ est surchargé et
peut se lire $M$}

\note{au milieu de la page, deux diagrammes en perspective, du type cube,
l'un sur les lettres $S$, $W$, $T$, $P$, $Q$, $M$, $R$, $N$, $(s,t)$, $t$,
$\emptyset$, l'autre sur les lettres correspondantes $V'$, $X'_{\eta}$,
$V'_{s}$, $V'_{s} \times \eta$, $X'_{\eta}$, $e$, $(s,t)$, $t$ ; plusieurs
régions hachurées, des flèches repassées. Ils ne sont pas reproduits}

\marginal{$D^{1} \times B \sim B$, $\ S^{1} \times B$ ; \struck{$H^{i}(B) \times H^{n-i}(B)$} $\ H^{i}(B)$}

\begin{tikzcd}[column sep=tiny, row sep=small, nodes={font=\scriptsize}]
 & H^{i-1}(R) \arrow[d] & & \\
H^{n-i}(P)^\vee \arrow[r] & H^{i}(W) \arrow[dl] \arrow[dr] & H^{n-i}(Q)^\vee \arrow[l] & \\
H^{i}(P) \arrow[r] & H^{i}(R) \arrow[d] & H^{i}(Q) \arrow[l] \arrow[r] & H^{n-i-1}(P)^\vee \\
 & H^{i+1}(W) & &
\end{tikzcd}

\note{l'exposant de $H^{n-i-1}(P)^\vee$ et celui de $H^{i+1}(W)$ sont écrits
avec un signe illisible devant le $1$ ; lectures douteuses}

\[
  H^{i}(Q) \times H^{n-i-1}(P) \longrightarrow H^{n-1}(R) \simeq \Lambda
\]
\[
  H^{n-i}(P)^\vee \longrightarrow H^{i}(P)
\]
\note{dans la première ligne, l'exposant $n-i-1$ est taché ; lecture douteuse}

\[
  \binom{5}{2} = \frac{5.4.3}{12.} = 10
\]
\note{tel quel sur la page}
\[
  \begin{array}{ccc}
    123 & 234 & 345 \\
    124 & 235 & \\
    125 & 245 & \\
    134 & & \\
    135 & & \\
    145 & &
  \end{array}
  \qquad 10 + 10 + 2
\]
\[
  \begin{array}{ccc}
    \emptyset & T \sim S & R \\
    5 + 5 & 2 + 2 & 2 \\
    3 & 3 + 3 & \\
    2 + 2 & 1 + 1 & \\
    1 + 2 & 1 + 1 & \\ \hline
    \boxed{20} & \boxed{14} & \boxed{2}
  \end{array}
  \qquad
  \begin{array}{c}
    M \sim N \\
    2 + 2 \\ \hline
    \boxed{4}
  \end{array}
\]

\marginal{40 suites exactes de coh. relative et
autant \uncertain{pour} l'homologie \ldots}\note{le nombre 40 est entouré}

En fait 80 suites exactes, mais

\page{19}

(a)

\begin{tikzcd}[column sep=small]
 & L^{\bullet\,\mathrm{triv}} = \mathbb{R}\Gamma(W_0)^{\mathrm{triv}} \arrow[dl, dashed] & \\
{[\mathbb{R}^\bullet]} \simeq \mathbb{R}(\Gamma, \pi)_{!}(\mathring{M}) \arrow[rr, dashed] & & K^\bullet = \mathbb{R}(\Gamma, \pi)(M) \arrow[ul, "\rho"']
\end{tikzcd}

\note{le sommet $\mathbb{R}^\bullet$ est mis entre crochets et entouré au
crayon ; l'isomorphisme avec $\mathbb{R}(\Gamma, \pi)_{!}(\mathring{M})$ est
écrit en dessous. La flèche de $L^{\bullet\,\mathrm{triv}}$ vers
$\mathbb{R}^\bullet$ est ondulée. Dans $\mathbb{R}\Gamma(W_0)$, deux signes
fortement biffés entre $\mathbb{R}$ et $\Gamma$, et un indice biffé ; de même
sous $\mathbb{R}(\Gamma,\pi)(M)$. Le long de $\rho$ : « compatible avec
cup »}

\marginal{on met dessus \uncertain{opér.} triviales de $\pi$}\note{relié par une
flèche à l'exposant « triv »}

dans catégorie dérivée de $\Lambda[\pi]$

(b) Appliquant $\mathbb{R}\Gamma^{\pi}$, on trouve
\note{au-dessus, relié à $\mathbb{R}\Gamma^{\pi}$ par un trait : « $\simeq$ (a) »}
\[
  L^{\bullet\pi} = \mathbb{R}\Gamma(W_0 \times S^{1}) \simeq
  \mathbb{R}\Gamma(W_0) \overset{\mathbf{L}}{\otimes}_{\Lambda} \mathbb{R}\Gamma(S^{1})
  \simeq \mathbb{R}\Gamma(W_0) \times \mathbb{R}\Gamma(W_0)(-1)
\]

\begin{tikzcd}[column sep=small]
 & L^{\bullet\pi} \arrow[dl, dashed] & \\
{[\mathbb{R}^{\bullet\pi}]} = \mathbb{R}\Gamma_{!}(\mathring{P}) \arrow[rr, dashed] & & K^{\bullet\pi} = \mathbb{R}\Gamma(P) = \mathbb{R}\Gamma(W - W_0) \arrow[ul, "{(\alpha,\ \beta)}"']
\end{tikzcd}

\note{$\mathbb{R}^{\bullet\pi}$ entre crochets et entouré au crayon ; sous
$\mathbb{R}\Gamma_{!}(\mathring{P})$ : « $= \mathbb{R}\Gamma_{!}(W - W_0)$ »}

\marginal{suite exacte de cohomologie relative pour
$(P, \partial P = S^{1} \times W_0)$}

\[
  \begin{aligned}
    \mathbb{R}\Gamma(P) \simeq \mathbb{R}\Gamma(W - W_0) &\xrightarrow[\ \alpha\ ]{\text{compatible avec cup}} \mathbb{R}\Gamma(W_0) \\
    \mathbb{R}\Gamma(P) \simeq \mathbb{R}\Gamma(W - W_0) &\xrightarrow{\ \beta\ } \mathbb{R}\Gamma(W_0)(-1)
  \end{aligned}
\]
\marginal{s'identifie dans la cohomologie au $\partial$ de la suite exacte de
cohomologie relative qui relie $H^{i}(W - W_0)$ et $H^{i-1}(W_0)$ \ldots}\note{sous
$\beta$ ; le décalage est écrit ici $(-1)$, entre parenthèses}

(c) Appliquant le foncteur oubli à (a), on trouve

\begin{tikzcd}[column sep=small]
 & L^\bullet = \mathbb{R}\Gamma(W_0) \arrow[dl, dashed] & \\
{[\mathbb{R}^{\omega}]} \simeq \mathbb{R}\Gamma_{!}(\mathring{M}) \arrow[rr, dashed] & & K^{\bullet\omega} = \mathbb{R}\Gamma(M) \arrow[ul, "\text{compatible avec cup}"']
\end{tikzcd}

\note{l'exposant, lu $\omega$, peut être un « w » ; $\mathbb{R}^{\omega}$ est
entre crochets et entouré au crayon}

(d)

\begin{tikzcd}[column sep=small]
 & \mathbb{R}\Gamma(P) \simeq \mathbb{R}\Gamma(W - W_0) \arrow[dl, "\beta"'] & \\
\mathbb{R}\Gamma(W_0)[-2] \arrow[rr] & & {[\mathbb{R}\Gamma_{!}(W)]} \arrow[ul, "\text{compatible avec cup}"']
\end{tikzcd}

\marginal{N.B. c'est le $\beta$ de (b)}\note{le long de la flèche de droite,
un début de mot est biffé}

[Suite exacte de cohomologie relative pour $(W, W - W_0)$, i.e. au choix pour
$(W, P)$] \ill{} \ill{} \ill{} $\mathbb{R}\Gamma(W)$ \ill{} \ill{} (a)

\end{document}
